\section{EXP-009 design (planned)}
\label{app:exp009-design}

This appendix expands the EXP-009 sketch in \S\ref{sec:followup} into the
operational details required to launch the experiment. The design is
committed in advance --- including the decision table mapping outcome
bands to claim revisions --- so the experiment cannot be re-interpreted
post-hoc.

\subsection{Goal}
Self-pretrain a vision encoder on a non-HAM10000 dermoscopy corpus, then
re-run the EXP-004 recipe on top with the linear predictor. The encoder
sees no HAM10000 image during pretraining, so the EXP-007-style positive
on \strongholdtwo\ --- if it appears --- can no longer be attributed to
HAM10000 image-level overlap and must reflect dermoscopy-domain transfer.

\subsection{Pretraining corpus}
Assemble from publicly available dermoscopy / clinical-skin sources, with
provenance audit on every component to confirm HAM10000 exclusion:
\begin{itemize}
\item ISIC archive components other than HAM10000 (BCN20000
      \citep{combalia2019bcn20000}, MSK-1 through MSK-5, UDA, Mednode,
      and similar legacy archives).
\item DermNet NZ and DermQuest-derived public images (clinical, not
      dermoscopic; included to broaden the visual distribution).
\item PAD-UFES-20 and 7-point checklist datasets (each cross-checked
      for HAM10000 overlap by lesion-ID and SHA-256 hash).
\end{itemize}
Provenance audit: every image must carry a source attribution and a
SHA-256 hash. The HAM10000 image-hash set is computed from the source
manifest and excluded by hash; any image whose source attribution is
ambiguous is excluded by default.

\subsection{Pretraining objective}
DINOv2 ViT-B/14 architecture (i.e.\ the same family as the EXP-004
backbone, so any test-$\auroc$ change is attributable to pretraining
data, not architecture), trained for $\sim$20 epochs with a JEPA-style
or masked-image-modelling objective. Hyperparameters target a
$\sim$$200{,}000$-step regime on a $\sim$$8\times$~A10G or
$1\times$~A100 budget. Pretrain-quality is checked by linear-probing on
a held-out HAM10000 split for diagnosis classification --- a sanity
check that the encoder has not overfit to its training distribution.

\subsection{Probe}
After pretraining, the encoder is frozen and the EXP-004 recipe runs on
top: linear predictor with identity warm-start, SGD, weight-decay-on-$(W-I)$,
$200$ epochs, the same lesion-ID splits, the same $1{,}000 + 1{,}000$
pairs per split, the same three-family augmentation pipeline, the same
seed. The only varied axis relative to EXP-004 is the pretraining
data; the only varied axis relative to EXP-007 is HAM10000 inclusion in
that pretraining data.

\subsection{Decision table}
Identical to Table~\ref{tab:exp009-decision} in \S\ref{sec:followup}.
Outcome bands and claim revisions:
\begin{itemize}
\item Test $\auroc \geq 0.85$: dermoscopy-domain transfer is sufficient
      under this scaffold. Headline survives; contamination caveat
      downgrades from ``unpartitioned'' to ``ruled out.''
\item Test $\auroc$ in $[0.50, 0.80]$: partial transfer; both factors
      contribute. Claim becomes ``dermoscopy-domain pretraining helps,
      HAM10000 image overlap helps more.''
\item Test $\auroc$ in $[0.30, 0.50]$: HAM10000 overlap drove EXP-007.
      Claim narrows to ``frozen-backbone JEPA unlocks when the backbone
      has seen the eval dataset, not when it has seen the eval
      domain.''
\end{itemize}

\subsection{Why not run before this preprint}
Three concrete costs gate the launch:
\begin{enumerate}[label=(\arabic*), leftmargin=2em]
\item \emph{Provenance audit.} Confirming HAM10000 exclusion on every
      ISIC component requires a per-image hash check against the
      HAM10000 manifest --- the audit is straightforward but is itself
      a multi-day engineering task before any GPU is spent.
\item \emph{Pretrain compute.} DINOv2 ViT-B/14 pretraining on
      $\sim$$200{,}000$ steps requires a non-trivial GPU budget that is
      not in the free-tier envelope used for the other nine runs.
\item \emph{Pretrain-quality validation.} A weak pretraining run that
      never learns dermoscopy structure would produce a low test
      $\auroc$ that the decision table would mis-attribute to (ii) or
      (iii) when the actual cause is pretrain failure. We need a
      held-out diagnostic to disambiguate.
\end{enumerate}
We share the methodology and the nine-run arc at this stage to invite
feedback on the proxy-task design and on whether the proposed partition
is the right one before committing to the build.
